\documentclass[10pt,a4paper]{amsart}
\usepackage{amsmath,amssymb,amsthm,graphicx}
% For commutative diagrams we use Xy-pic
\usepackage[all]{xy}

% We use 2cell for 2-commutative diagrams.
\xyoption{2cell}
% Theorem environments.
%
\theoremstyle{plain}
\newtheorem{theorem}[subsection]{Theorem}
\newtheorem{proposition}[subsection]{Proposition}
\newtheorem{lemma}[subsection]{Lemma}

\theoremstyle{definition}
\newtheorem{definition}[subsection]{Definition}
\newtheorem{example}[subsection]{Example}
\newtheorem{exercise}[subsection]{Exercise}
\newtheorem{situation}[subsection]{Situation}

\theoremstyle{remark}
\newtheorem{remark}[subsection]{Remark}
\newtheorem{remarks}[subsection]{Remarks}

\numberwithin{equation}{subsection}

% Macros
%
\def\lim{\mathop{\mathrm{lim}}\nolimits}
\def\colim{\mathop{\mathrm{colim}}\nolimits}
\def\Spec{\mathop{\mathrm{Spec}}}
\def\Hom{\mathop{\mathrm{Hom}}\nolimits}
\def\Ext{\mathop{\mathrm{Ext}}\nolimits}
\def\SheafHom{\mathop{\mathcal{H}\!\mathit{om}}\nolimits}
\def\SheafExt{\mathop{\mathcal{E}\!\mathit{xt}}\nolimits}
\def\Sch{\mathit{Sch}}
\def\Mor{\mathop{\mathrm{Mor}}\nolimits}
\def\Ob{\mathop{\mathrm{Ob}}\nolimits}
\def\Sh{\mathop{\mathit{Sh}}\nolimits}
\def\NL{\mathop{N\!L}\nolimits}
\def\CH{\mathop{\mathrm{CH}}\nolimits}
\def\proetale{{pro\text{-}\acute{e}tale}}
\def\etale{{\acute{e}tale}}
\def\QCoh{\mathit{QCoh}}
\def\Ker{\mathop{\mathrm{Ker}}}
\def\Im{\mathop{\mathrm{Im}}}
\def\Coker{\mathop{\mathrm{Coker}}}
\def\Coim{\mathop{\mathrm{Coim}}}

% Boxtimes
%
\DeclareMathSymbol{\boxtimes}{\mathbin}{AMSa}{"02}

%
% Macros for moduli stacks/spaces
%
\def\QCohstack{\mathcal{QC}\!\mathit{oh}}
\def\Cohstack{\mathcal{C}\!\mathit{oh}}
\def\Spacesstack{\mathcal{S}\!\mathit{paces}}
\def\Quotfunctor{\mathrm{Quot}}
\def\Hilbfunctor{\mathrm{Hilb}}
\def\Curvesstack{\mathcal{C}\!\mathit{urves}}
\def\Polarizedstack{\mathcal{P}\!\mathit{olarized}}
\def\Complexesstack{\mathcal{C}\!\mathit{omplexes}}
% \Pic is the operator that assigns to X its picard group, usage \Pic(X)
% \Picardstack_{X/B} denotes the Picard stack of X over B
% \Picardfunctor_{X/B} denotes the Picard functor of X over B
\def\Pic{\mathop{\mathrm{Pic}}\nolimits}
\def\Picardstack{\mathcal{P}\!\mathit{ic}}
\def\Picardfunctor{\mathrm{Pic}}
\def\Deformationcategory{\mathcal{D}\!\mathit{ef}}
\newenvironment{reference}{}{}
\newenvironment{slogan}{}{}
\newenvironment{history}{}{}
\usepackage{cleveref}
\crefname{originallemma}{original lemma}{original lemmas}
\Crefname{originallemma}{Original lemma}{Original lemmas}
\makeatletter
\newlabel{lemma-genus-nonnegative}{{0C7C}{source}{}{sourcepointer-1}{}}
\newlabel{lemma-genus-nonnegative@cref}{{[originallemma][0][]0C7C}{[1][0][]source}{}{}{}}
\makeatother
\crefname{originallemma}{original lemma}{original lemmas}
\Crefname{originallemma}{Original lemma}{Original lemmas}
\makeatletter
\newlabel{item-one}{{0C8U}{source}{}{sourcepointer-2}{}}
\newlabel{item-one@cref}{{[originallemma][0][]0C8U}{[1][0][]source}{}{}{}}
\makeatother
\crefname{originallemma}{original lemma}{original lemmas}
\Crefname{originallemma}{Original lemma}{Original lemmas}
\makeatletter
\newlabel{lemma-genus-one}{{0C8T}{source}{}{sourcepointer-3}{}}
\newlabel{lemma-genus-one@cref}{{[originallemma][0][]0C8T}{[1][0][]source}{}{}{}}
\makeatother
\crefname{originallemma}{original lemma}{original lemmas}
\Crefname{originallemma}{Original lemma}{Original lemmas}
\makeatletter
\newlabel{lemma-recurring-symmetric-real}{{0C5X}{source}{}{sourcepointer-4}{}}
\newlabel{lemma-recurring-symmetric-real@cref}{{[originallemma][0][]0C5X}{[1][0][]source}{}{}{}}
\makeatother
\crefname{originaldefinition}{original definition}{original definitions}
\Crefname{originaldefinition}{Original definition}{Original definitions}
\makeatletter
\newlabel{varieties-definition-algebraic-scheme}{{06LG}{source}{}{sourcepointer-5}{}}
\newlabel{varieties-definition-algebraic-scheme@cref}{{[originaldefinition][0][]06LG}{[1][0][]source}{}{}{}}
\makeatother
\crefname{originalsection}{original section}{original sections}
\Crefname{originalsection}{Original section}{Original sections}
\makeatletter
\newlabel{resolve-section-minus-one}{{0C2I}{source}{}{sourcepointer-6}{}}
\newlabel{resolve-section-minus-one@cref}{{[originalsection][0][]0C2I}{[1][0][]source}{}{}{}}
\makeatother
\crefname{originallemma}{original lemma}{original lemmas}
\Crefname{originallemma}{Original lemma}{Original lemmas}
\makeatletter
\newlabel{lemma-divisor-special-fiber}{{0C5Z}{source}{}{sourcepointer-7}{}}
\newlabel{lemma-divisor-special-fiber@cref}{{[originallemma][0][]0C5Z}{[1][0][]source}{}{}{}}
\makeatother
\crefname{originallemma}{original lemma}{original lemmas}
\Crefname{originallemma}{Original lemma}{Original lemmas}
\makeatletter
\newlabel{lemma-genus-formula}{{0CA3}{source}{}{sourcepointer-8}{}}
\newlabel{lemma-genus-formula@cref}{{[originallemma][0][]0CA3}{[1][0][]source}{}{}{}}
\makeatother
\crefname{originalequation}{original equation}{original equations}
\Crefname{originalequation}{Original equation}{Original equations}
\makeatletter
\newlabel{equation-form}{{0C65}{source}{}{sourcepointer-9}{}}
\newlabel{equation-form@cref}{{[originalequation][0][]0C65}{[1][0][]source}{}{}{}}
\makeatother
\crefname{originalsection}{original section}{original sections}
\Crefname{originalsection}{Original section}{Original sections}
\makeatletter
\newlabel{lemma-intersection-pairing}{{0C64}{source}{}{sourcepointer-10}{}}
\newlabel{lemma-intersection-pairing@cref}{{[originalsection][0][]0C64}{[1][0][]source}{}{}{}}
\makeatother
\crefname{originallemma}{original lemma}{original lemmas}
\Crefname{originallemma}{Original lemma}{Original lemmas}
\makeatletter
\newlabel{lemma-properties-form}{{0C66}{source}{}{sourcepointer-11}{}}
\newlabel{lemma-properties-form@cref}{{[originallemma][0][]0C66}{[1][0][]source}{}{}{}}
\makeatother
\crefname{originallemma}{original lemma}{original lemmas}
\Crefname{originallemma}{Original lemma}{Original lemmas}
\makeatletter
\newlabel{varieties-lemma-divisible}{{0C6P}{source}{}{sourcepointer-12}{}}
\newlabel{varieties-lemma-divisible@cref}{{[originallemma][0][]0C6P}{[1][0][]source}{}{}{}}
\makeatother
\crefname{originalproposition}{original proposition}{original propositions}
\Crefname{originalproposition}{Original proposition}{Original propositions}
\makeatletter
\newlabel{curves-proposition-projective-line}{{0C6U}{source}{}{sourcepointer-13}{}}
\newlabel{curves-proposition-projective-line@cref}{{[originalproposition][0][]0C6U}{[1][0][]source}{}{}{}}
\makeatother
\crefname{originallemma}{original lemma}{original lemmas}
\Crefname{originallemma}{Original lemma}{Original lemmas}
\makeatletter
\newlabel{resolve-lemma-birational-regular-surfaces}{{0C5S}{source}{}{sourcepointer-14}{}}
\newlabel{resolve-lemma-birational-regular-surfaces@cref}{{[originallemma][0][]0C5S}{[1][0][]source}{}{}{}}
\makeatother
\crefname{originallemma}{original lemma}{original lemmas}
\Crefname{originallemma}{Original lemma}{Original lemmas}
\makeatletter
\newlabel{varieties-lemma-modification-normal-iso-over-codimension-1}{{0BFP}{source}{}{sourcepointer-15}{}}
\newlabel{varieties-lemma-modification-normal-iso-over-codimension-1@cref}{{[originallemma][0][]0BFP}{[1][0][]source}{}{}{}}
\makeatother
\crefname{originallemma}{original lemma}{original lemmas}
\Crefname{originallemma}{Original lemma}{Original lemmas}
\makeatletter
\newlabel{resolve-lemma-proper-birational-regular-surfaces}{{0C5R}{source}{}{sourcepointer-16}{}}
\newlabel{resolve-lemma-proper-birational-regular-surfaces@cref}{{[originallemma][0][]0C5R}{[1][0][]source}{}{}{}}
\makeatother
\crefname{originallemma}{original lemma}{original lemmas}
\Crefname{originallemma}{Original lemma}{Original lemmas}
\makeatletter
\newlabel{resolve-lemma-contract-when-quasi-projective}{{0C2N}{source}{}{sourcepointer-17}{}}
\newlabel{resolve-lemma-contract-when-quasi-projective@cref}{{[originallemma][0][]0C2N}{[1][0][]source}{}{}{}}
\makeatother
\crefname{originallemma}{original lemma}{original lemmas}
\Crefname{originallemma}{Original lemma}{Original lemmas}
\makeatletter
\newlabel{resolve-lemma-factor-through-contraction}{{0C5J}{source}{}{sourcepointer-18}{}}
\newlabel{resolve-lemma-factor-through-contraction@cref}{{[originallemma][0][]0C5J}{[1][0][]source}{}{}{}}
\makeatother
\crefname{originallemma}{original lemma}{original lemmas}
\Crefname{originallemma}{Original lemma}{Original lemmas}
\makeatletter
\newlabel{algebra-lemma-regular-ring-CM}{{00NQ}{source}{}{sourcepointer-19}{}}
\newlabel{algebra-lemma-regular-ring-CM@cref}{{[originallemma][0][]00NQ}{[1][0][]source}{}{}{}}
\makeatother
\crefname{originallemma}{original lemma}{original lemmas}
\Crefname{originallemma}{Original lemma}{Original lemmas}
\makeatletter
\newlabel{varieties-lemma-regular-point-on-curve}{{0B8Y}{source}{}{sourcepointer-20}{}}
\newlabel{varieties-lemma-regular-point-on-curve@cref}{{[originallemma][0][]0B8Y}{[1][0][]source}{}{}{}}
\makeatother
\crefname{originallemma}{original lemma}{original lemmas}
\Crefname{originallemma}{Original lemma}{Original lemmas}
\makeatletter
\newlabel{divisors-lemma-strict-transform}{{080E}{source}{}{sourcepointer-21}{}}
\newlabel{divisors-lemma-strict-transform@cref}{{[originallemma][0][]080E}{[1][0][]source}{}{}{}}
\makeatother
\crefname{originallemma}{original lemma}{original lemmas}
\Crefname{originallemma}{Original lemma}{Original lemmas}
\makeatletter
\newlabel{divisors-lemma-blow-up-effective-Cartier-divisor}{{0807}{source}{}{sourcepointer-22}{}}
\newlabel{divisors-lemma-blow-up-effective-Cartier-divisor@cref}{{[originallemma][0][]0807}{[1][0][]source}{}{}{}}
\makeatother
\title{Nonunique genus-zero minimal models have projective-line special fibres}
\author{The Stacks Project Authors}
\date{stacks-project:models.tex:local-snapshot-20261011; GNU FDL 1.2 or later}
\begin{document}
\maketitle
\noindent\textit{Editorial note.} Original author theorem and proof are retained. Only this target is counted; local supporting statements are not separate new entries.\par\medskip
\noindent\textit{Allowed standard prerequisites (editorial).}\begin{itemize}
\item DVR models of smooth projective curves, common resolutions of regular surfaces, contraction of exceptional curves of the first kind, strict transforms and regular local rings. Original numerical-type and minimality support retained.. These are the named standard inputs; the selected main construction and its local nonroutine support remain in full.
\end{itemize}\medskip
\ifdefined\hypertarget\hypertarget{sourcepointer-1}{}\fi
\noindent\textit{Source reference.} Original lemma 0C7C: focused-seeds/AGP/sources/stacks-snapshot-20261011/tags.txt, source lines 15327--15327. Stable Stacks locator for declared standard input or surrounding context. No selected local proof is replaced by this pointer.\par\medskip
\ifdefined\hypertarget\hypertarget{sourcepointer-2}{}\fi
\noindent\textit{Source reference.} Original lemma 0C8U: focused-seeds/AGP/sources/stacks-snapshot-20261011/tags.txt, source lines 15379--15379. Stable Stacks locator for declared standard input or surrounding context. No selected local proof is replaced by this pointer.\par\medskip
\ifdefined\hypertarget\hypertarget{sourcepointer-3}{}\fi
\noindent\textit{Source reference.} Original lemma 0C8T: focused-seeds/AGP/sources/stacks-snapshot-20261011/tags.txt, source lines 15378--15378. Stable Stacks locator for declared standard input or surrounding context. No selected local proof is replaced by this pointer.\par\medskip
\ifdefined\hypertarget\hypertarget{sourcepointer-4}{}\fi
\noindent\textit{Source reference.} Original lemma 0C5X: focused-seeds/AGP/sources/stacks-snapshot-20261011/tags.txt, source lines 15277--15277. Stable Stacks locator for declared standard input or surrounding context. No selected local proof is replaced by this pointer.\par\medskip
\ifdefined\hypertarget\hypertarget{sourcepointer-5}{}\fi
\noindent\textit{Source reference.} Original definition 06LG: focused-seeds/AGP/sources/stacks-snapshot-20261011/tags.txt, source lines 8465--8465. Stable Stacks locator for declared standard input or surrounding context. No selected local proof is replaced by this pointer.\par\medskip
\ifdefined\hypertarget\hypertarget{sourcepointer-6}{}\fi
\noindent\textit{Source reference.} Original section 0C2I: focused-seeds/AGP/sources/stacks-snapshot-20261011/tags.txt, source lines 15158--15158. Stable Stacks locator for declared standard input or surrounding context. No selected local proof is replaced by this pointer.\par\medskip
\ifdefined\hypertarget\hypertarget{sourcepointer-7}{}\fi
\noindent\textit{Source reference.} Original lemma 0C5Z: focused-seeds/AGP/sources/stacks-snapshot-20261011/tags.txt, source lines 15279--15279. Stable Stacks locator for declared standard input or surrounding context. No selected local proof is replaced by this pointer.\par\medskip
\ifdefined\hypertarget\hypertarget{sourcepointer-8}{}\fi
\noindent\textit{Source reference.} Original lemma 0CA3: focused-seeds/AGP/sources/stacks-snapshot-20261011/tags.txt, source lines 15423--15423. Stable Stacks locator for declared standard input or surrounding context. No selected local proof is replaced by this pointer.\par\medskip
\ifdefined\hypertarget\hypertarget{sourcepointer-9}{}\fi
\noindent\textit{Source reference.} Original equation 0C65: focused-seeds/AGP/sources/stacks-snapshot-20261011/tags.txt, source lines 15285--15285. Stable Stacks locator for declared standard input or surrounding context. No selected local proof is replaced by this pointer.\par\medskip
\ifdefined\hypertarget\hypertarget{sourcepointer-10}{}\fi
\noindent\textit{Source reference.} Original section 0C64: focused-seeds/AGP/sources/stacks-snapshot-20261011/tags.txt, source lines 15284--15284. Stable Stacks locator for declared standard input or surrounding context. No selected local proof is replaced by this pointer.\par\medskip
\ifdefined\hypertarget\hypertarget{sourcepointer-11}{}\fi
\noindent\textit{Source reference.} Original lemma 0C66: focused-seeds/AGP/sources/stacks-snapshot-20261011/tags.txt, source lines 15286--15286. Stable Stacks locator for declared standard input or surrounding context. No selected local proof is replaced by this pointer.\par\medskip
\ifdefined\hypertarget\hypertarget{sourcepointer-12}{}\fi
\noindent\textit{Source reference.} Original lemma 0C6P: focused-seeds/AGP/sources/stacks-snapshot-20261011/tags.txt, source lines 15304--15304. Stable Stacks locator for declared standard input or surrounding context. No selected local proof is replaced by this pointer.\par\medskip
\ifdefined\hypertarget\hypertarget{sourcepointer-13}{}\fi
\noindent\textit{Source reference.} Original proposition 0C6U: focused-seeds/AGP/sources/stacks-snapshot-20261011/tags.txt, source lines 15309--15309. Stable Stacks locator for declared standard input or surrounding context. No selected local proof is replaced by this pointer.\par\medskip
\ifdefined\hypertarget\hypertarget{sourcepointer-14}{}\fi
\noindent\textit{Source reference.} Original lemma 0C5S: focused-seeds/AGP/sources/stacks-snapshot-20261011/tags.txt, source lines 15272--15272. Stable Stacks locator for declared standard input or surrounding context. No selected local proof is replaced by this pointer.\par\medskip
\ifdefined\hypertarget\hypertarget{sourcepointer-15}{}\fi
\noindent\textit{Source reference.} Original lemma 0BFP: focused-seeds/AGP/sources/stacks-snapshot-20261011/tags.txt, source lines 14392--14392. Stable Stacks locator for declared standard input or surrounding context. No selected local proof is replaced by this pointer.\par\medskip
\ifdefined\hypertarget\hypertarget{sourcepointer-16}{}\fi
\noindent\textit{Source reference.} Original lemma 0C5R: focused-seeds/AGP/sources/stacks-snapshot-20261011/tags.txt, source lines 15271--15271. Stable Stacks locator for declared standard input or surrounding context. No selected local proof is replaced by this pointer.\par\medskip
\ifdefined\hypertarget\hypertarget{sourcepointer-17}{}\fi
\noindent\textit{Source reference.} Original lemma 0C2N: focused-seeds/AGP/sources/stacks-snapshot-20261011/tags.txt, source lines 15163--15163. Stable Stacks locator for declared standard input or surrounding context. No selected local proof is replaced by this pointer.\par\medskip
\ifdefined\hypertarget\hypertarget{sourcepointer-18}{}\fi
\noindent\textit{Source reference.} Original lemma 0C5J: focused-seeds/AGP/sources/stacks-snapshot-20261011/tags.txt, source lines 15264--15264. Stable Stacks locator for declared standard input or surrounding context. No selected local proof is replaced by this pointer.\par\medskip
\ifdefined\hypertarget\hypertarget{sourcepointer-19}{}\fi
\noindent\textit{Source reference.} Original lemma 00NQ: focused-seeds/AGP/sources/stacks-snapshot-20261011/tags.txt, source lines 868--868. Stable Stacks locator for declared standard input or surrounding context. No selected local proof is replaced by this pointer.\par\medskip
\ifdefined\hypertarget\hypertarget{sourcepointer-20}{}\fi
\noindent\textit{Source reference.} Original lemma 0B8Y: focused-seeds/AGP/sources/stacks-snapshot-20261011/tags.txt, source lines 14156--14156. Stable Stacks locator for declared standard input or surrounding context. No selected local proof is replaced by this pointer.\par\medskip
\ifdefined\hypertarget\hypertarget{sourcepointer-21}{}\fi
\noindent\textit{Source reference.} Original lemma 080E: focused-seeds/AGP/sources/stacks-snapshot-20261011/tags.txt, source lines 10179--10179. Stable Stacks locator for declared standard input or surrounding context. No selected local proof is replaced by this pointer.\par\medskip
\ifdefined\hypertarget\hypertarget{sourcepointer-22}{}\fi
\noindent\textit{Source reference.} Original lemma 0807: focused-seeds/AGP/sources/stacks-snapshot-20261011/tags.txt, source lines 10172--10172. Stable Stacks locator for declared standard input or surrounding context. No selected local proof is replaced by this pointer.\par\medskip
\section{Numerical types}
\label{section-numerical-types}

\noindent
Part of the arguments will involve the combinatorics of the following
data structures.

\begin{definition}
\label{definition-type}
A {\it numerical type} $T$ is given by
$$
n, m_i, a_{ij}, w_i, g_i
$$
where $n \geq 1$ is an integer and $m_i$, $a_{ij}$, $w_i$, $g_i$
are integers for $1 \leq i, j \leq n$ subject to the following conditions
\begin{enumerate}
\item $m_i > 0$, $w_i > 0$, $g_i \geq 0$,
\item the matrix $A = (a_{ij})$ is symmetric and $a_{ij} \geq 0$
for $i \not = j$,
\item there is no proper nonempty subset $I \subset \{1, \ldots, n\}$
such that $a_{ij} = 0$ for $i \in I$, $j \not \in I$,
\item for each $i$ we have $\sum_j a_{ij}m_j = 0$, and
\item $w_i | a_{ij}$.
\end{enumerate}
\end{definition}

\noindent
This is obviously a somewhat annoying type of structure to work with,
but it is exactly what shows up in special fibres of proper regular
models of smooth geometrically connected curves.
Of course we only care about these types up to reordering the indices.

\begin{definition}
\label{definition-type-equivalent}
We say two numerical types $n, m_i, a_{ij}, w_i, g_i$ and
$n', m'_i, a'_{ij}, w'_i, g'_i$ are {\it equivalent types} if
there exists a permutation $\sigma$ of $\{1, \ldots, n\}$
such that $m_i = m'_{\sigma(i)}$, $a_{ij} = a'_{\sigma(i)\sigma(j)}$,
$w_i = w'_{\sigma(i)}$, and $g_i = g'_{\sigma(i)}$.
\end{definition}

\noindent
A numerical type has a genus.

\begin{lemma}
\label{lemma-genus}
Let $n, m_i, a_{ij}, w_i, g_i$ be a numerical type. Then the expression
$$
g = 1 + \sum m_i(w_i(g_i - 1) - \frac{1}{2} a_{ii})
$$
is an integer.
\end{lemma}

\begin{proof}
To prove $g$ is an integer we have to show that $\sum a_{ii}m_i$ is even.
This we can see by computing modulo $2$ as follows
\begin{align*}
\sum\nolimits_i a_{ii} m_i 
& \equiv
\sum\nolimits_{i,\ m_i\text{ odd}} a_{ii}m_i \\
& \equiv
\sum\nolimits_{i,\ m_i\text{ odd}} \sum\nolimits_{j \not = i} a_{ij}m_j \\
& \equiv
\sum\nolimits_{i,\ m_i\text{ odd}}
\sum\nolimits_{j \not = i,\ m_j\text{ odd}} a_{ij}m_j \\
& \equiv
\sum\nolimits_{i < j,\ m_i\text{ and }m_j\text{ odd}} a_{ij}(m_i + m_j) \\
& \equiv
0
\end{align*}
where we have used that $a_{ij} = a_{ji}$ and that $\sum_j a_{ij}m_j = 0$
for all $i$.
\end{proof}

\begin{definition}
\label{definition-genus}
We say $n, m_i, a_{ij}, w_i, g_i$ is a {\it numerical type of genus $g$}
if $g = 1 + \sum m_i(w_i(g_i - 1) - \frac{1}{2} a_{ii})$ is the integer
from Lemma \ref{lemma-genus}.
\end{definition}

\noindent
We will prove below (Lemma \ref{lemma-genus-nonnegative}) that the genus
is almost always $\geq 0$. But you can have numerical types with
negative genus.

\begin{lemma}
\label{lemma-irreducible}
Let $n, m_i, a_{ij}, w_i, g_i$ be a numerical type of genus $g$.
If $n = 1$, then $a_{11} = 0$ and $g = 1 + m_1w_1(g_1 - 1)$.
Moreover, we can classify all such numerical types as follows
\begin{enumerate}
\item If $g < 0$, then $g_1 = 0$ and there are finitely many possible
numerical types of genus $g$ with $n = 1$ corresponding to factorizations
$m_1w_1 = 1 - g$.
\item If $g = 0$, then $m_1 = 1$, $w_1 = 1$, $g_1 = 0$
as in Lemma \ref{lemma-genus-zero}.
\item If $g = 1$, then we conclude $g_1 = 1$ but $m_1, w_1$ can be arbitrary
positive integers; this is case
(\ref{item-one}) of Lemma \ref{lemma-genus-one}.
\item If $g > 1$, then $g_1 > 1$ and there are finitely many possible
numerical types of genus $g$ with $n = 1$ corresponding to
factorizations $m_1w_1(g_1 - 1) = g - 1$.
\end{enumerate}
\end{lemma}

\begin{proof}
The lemma proves itself.
\end{proof}

\begin{lemma}
\label{lemma-diagonal-negative}
Let $n, m_i, a_{ij}, w_i, g_i$ be a numerical type of genus $g$.
If $n > 1$, then $a_{ii} < 0$ for all $i$.
\end{lemma}

\begin{proof}
Lemma \ref{lemma-recurring-symmetric-real} applies to the matrix $A$.
\end{proof}

\begin{lemma}
\label{lemma-minus-one}
Let $n, m_i, a_{ij}, w_i, g_i$ be a numerical type of genus $g$.
Assume $n > 1$. If $i$ is such that the contribution
$m_i(w_i(g_i - 1) - \frac{1}{2} a_{ii})$
to the genus $g$ is $< 0$, then $g_i = 0$ and $a_{ii} = -w_i$.
\end{lemma}

\begin{proof}
Follows immediately from Lemma \ref{lemma-diagonal-negative} and
$w_i > 0$, $g_i \geq 0$, and $w_i | a_{ii}$.
\end{proof}

\begin{definition}
\label{definition-type-minus-one}
Let $n, m_i, a_{ij}, w_i, g_i$ be a numerical type.
We say $i$ is a {\it $(-1)$-index} if $g_i = 0$ and $a_{ii} = -w_i$.
\end{definition}

\noindent
We can ``contract'' $(-1)$-indices.

\begin{lemma}
\label{lemma-contract}
Let $n, m_i, a_{ij}, w_i, g_i$ be a numerical type $T$.
Assume $n$ is a $(-1)$-index. Then there is a numerical
type $T'$ given by $n', m'_i, a'_{ij}, w'_i, g'_i$ with
\begin{enumerate}
\item $n' = n - 1$,
\item $m'_i = m_i$,
\item $a'_{ij} = a_{ij} - a_{in}a_{jn}/a_{nn}$,
\item $w'_i = w_i/2$ if $a_{in}/w_n$ even and $a_{in}/w_i$ odd
and $w'_i = w_i$ else,
\item $g'_i =
\frac{w_i}{w'_i}(g_i - 1) + 1 + \frac{a_{in}^2 - w_na_{in}}{2w'_iw_n}$.
\end{enumerate}
Moreover, we have $g = g'$.
\end{lemma}

\begin{proof}
Observe that $n > 1$ for example by Lemma \ref{lemma-irreducible}
and hence $n' \geq 1$. We check conditions (1) -- (5) of
Definition \ref{definition-type} for $n', m'_i, a'_{ij}, w'_i, g'_i$.

\medskip\noindent
Condition (1) is immediate.

\medskip\noindent
Condition (2). Symmetry of $A' = (a'_{ij})$ is immediate
and since $a_{nn} < 0$ by Lemma \ref{lemma-diagonal-negative}
we see that $a'_{ij} \geq a_{ij} \geq 0$ if $i \not = j$.

\medskip\noindent
Condition (3). Suppose that $I \subset \{1, \ldots, n - 1\}$ such that
$a'_{ii'} = 0$ for $i \in I$ and $i' \in \{1, \ldots, n - 1\} \setminus I$.
Then we see that for each $i \in I$ and $i' \in I'$ we have
$a_{in}a_{i'n} = 0$. Thus either $a_{in} = 0$ for all $i \in I$ and
$I \subset \{1, \ldots, n\}$ is a contradiction for property (3) for $T$,
or $a_{i'n} = 0$ for all $i' \in \{1, \ldots, n - 1\} \setminus I$
and $I \cup \{n\} \subset \{1, \ldots, n\}$ is a contradiction for
property (3) of $T$. Hence (3) holds for $T'$.

\medskip\noindent
Condition (4). We compute
$$
\sum\nolimits_{j = 1}^{n - 1} a'_{ij}m_j =
\sum\nolimits_{j = 1}^{n - 1}
(a_{ij}m_j - \frac{a_{in}a_{jn}m_j}{a_{nn}}) =
- a_{in}m_n - \frac{a_{in}}{a_{nn}}(-a_{nn}m_n) = 0
$$
as desired.

\medskip\noindent
Condition (5). We have to show that $w'_i$ divides $a_{in}a_{jn}/a_{nn}$.
This is clear because $a_{nn} = -w_n$ and $w_n | a_{jn}$ and $w_i | a_{in}$.

\medskip\noindent
To show that $g = g'$ we first write
\begin{align*}
g
& =
1 + \sum\nolimits_{i = 1}^n m_i(w_i(g_i - 1) - \frac{1}{2}a_{ii}) \\
& =
1 + \sum\nolimits_{i = 1}^{n - 1} m_i(w_i(g_i - 1) - \frac{1}{2}a_{ii})
-\frac{1}{2}m_nw_n \\
& =
1 +  \sum\nolimits_{i = 1}^{n - 1}
m_i(w_i(g_i - 1) - \frac{1}{2}a_{ii} - \frac{1}{2}a_{in})
\end{align*}
Comparing with the expression for $g'$ we see that it suffices if
$$
w'_i(g'_i - 1) - \frac{1}{2}a'_{ii} =
w_i(g_i - 1) - \frac{1}{2}a_{in} - \frac{1}{2}a_{ii}
$$
for $i \leq n - 1$. In other words, we have
$$
g'_i = \frac{2w_i(g_i - 1) - a_{in} - a_{ii} + a'_{ii} + 2w'_i}{2w'_i} =
\frac{w_i}{w'_i}(g_i - 1) + 1 + \frac{a_{in}^2 - w_na_{in}}{2w'_iw_n}
$$
It is elementary to check that this is an integer $\geq 0$
if we choose $w'_i$ as in (4).
\end{proof}

\begin{lemma}
\label{lemma-top-genus}
Let $n, m_i, a_{ij}, w_i, g_i$ be a numerical type.
Let $e$ be the number of pairs $(i, j)$ with $i < j$ and $a_{ij} > 0$.
Then the expression $g_{top} = 1 - n + e$ is $\geq 0$.
\end{lemma}

\begin{proof}
If not, then $e < n - 1$ which means there exists an $i$ such that
$a_{ij} = 0$ for all $j \not = i$. This contradicts assumption
(3) of Definition \ref{definition-type}.
\end{proof}

\begin{definition}
\label{definition-top-genus}
Let $n, m_i, a_{ij}, w_i, g_i$ be a numerical type $T$. The
{\it topological genus of $T$} is the nonnegative integer
$g_{top} = 1 - n + e$ from Lemma \ref{lemma-top-genus}.
\end{definition}

\noindent
We want to bound the genus by the topological genus. However, this
will not always be the case, for example for numerical types
with $n = 1$ as in Lemma \ref{lemma-irreducible}. But it will
be true for minimal numerical types which are defined as follows.

\begin{definition}
\label{definition-type-minimal}
We say the numerical type $n, m_i, a_{ij}, w_i, g_i$ of genus $g$
is {\it minimal} if there does not exist an $i$
with $g_i = 0$ and $a_{ii} = -w_i$, in other words, if there
does not exist a $(-1)$-index.
\end{definition}

\noindent
We will prove that the genus $g$ of a minimal type with $n > 1$
is greater than or equal to $\max(1, g_{top})$.

\begin{lemma}
\label{lemma-non-irreducible-minimal-type-genus-at-least-one}
If $n, m_i, a_{ij}, w_i, g_i$ is a minimal numerical type
with $n > 1$, then $g \geq 1$.
\end{lemma}

\begin{proof}
This is true because $g = 1 + \sum \Phi_i$ with
$\Phi_i = m_i(w_i(g_i - 1) - \frac{1}{2} a_{ii})$ nonnegative
by Lemma \ref{lemma-minus-one} and the definition of minimal types.
\end{proof}

\begin{lemma}[Genus zero]
\label{lemma-genus-zero}
The only minimal numerical type of genus zero is
$n = 1$, $m_1 = 1$, $a_{11} = 0$, $w_1 = 1$, $g_1 = 0$.
\end{lemma}

\begin{proof}
Follows from Lemmas \ref{lemma-non-irreducible-minimal-type-genus-at-least-one}
and \ref{lemma-irreducible}.
\end{proof}

\section{Models}
\label{section-models}

\noindent
In this chapter $R$ will be a discrete valuation ring and $K$ will
be its fraction field. If needed we will denote $\pi \in R$ a
uniformizer and $k = R/(\pi)$ its residue field.

\medskip\noindent
Let $V$ be an algebraic $K$-scheme
(Varieties, Definition \ref{varieties-definition-algebraic-scheme}).
A {\it model} for $V$ will
mean a flat finite type\footnote{Occasionally it is useful to
allow models to be locally of finite type over $R$, but we'll
cross that bridge when we come to it.}
morphism $X \to \Spec(R)$ endowed with
an isomorphism $V \to X_K = X \times_{\Spec(R)} \Spec(K)$. We often
will identify $V$ and the generic fibre $X_K$ of $X$ and
just write $V = X_K$.
The special fibre is $X_k = X \times_{\Spec(R)} \Spec(k)$.
A {\it morphism of models $X \to X'$ for $V$}
is a morphism $X \to X'$ of schemes over $R$ which induces
the identity on $V$.

\medskip\noindent
We will say {\it $X$ is a proper model of $V$} if $X$ is a model
of $V$ and the structure morphism $X \to \Spec(R)$ is proper.
Similarly for separated models, smooth models, and add more here.
We will say {\it $X$ is a regular model of $V$} if $X$ is a model
of $V$ and $X$ is a regular scheme.
Similarly for normal models, reduced models, and add more here.

\begin{definition}
\label{definition-minimal-model}
Let $C$ be a smooth projective curve over $K$ with
$H^0(C, \mathcal{O}_C) = K$. A {\it minimal model}
will be a regular, proper model $X$ for $C$ such that
$X$ does not contain an exceptional curve of the first kind
(Resolution of Surfaces, Section \ref{resolve-section-minus-one}).
\end{definition}

\begin{situation}
\label{situation-regular-model}
Let $R$ be a discrete valuation ring with fraction field $K$,
residue field $k$, and uniformizer $\pi$.
Let $C$ be a smooth projective curve over $K$ with $H^0(C, \mathcal{O}_C) = K$.
Let $X$ be a regular proper model of $C$.
Let $C_1, \ldots, C_n$ be the irreducible components of the special
fibre $X_k$. Write $X_k = \sum m_i C_i$ as in
Lemma \ref{lemma-divisor-special-fiber}.
\end{situation}


\begin{lemma}
\label{lemma-numerical-type-of-model}
In Situation \ref{situation-regular-model} with
$\kappa_i = H^0(C_i, \mathcal{O}_{C_i})$ and $g_i$ the genus of $C_i$
the data
$$
n, m_i, (C_i \cdot C_j), [\kappa_i : k], g_i
$$
is a numerical type of genus equal to the genus of $C$.
\end{lemma}

\begin{proof}
(In the proof of Lemma \ref{lemma-genus-formula}
we have seen that the quantities
used in the statement of the lemma are well defined.)
We have to verify the conditions (1) -- (5) of
Definition \ref{definition-type}.

\medskip\noindent
Condition (1) is immediate.

\medskip\noindent
Condition (2). Symmetry of the matrix $(C_i \cdot C_j)$ follows from
Equation (\ref{equation-form}) and
Lemma \ref{lemma-intersection-pairing}.
Nonnegativity of $(C_i \cdot C_j)$ for $i \not = j$
is part (1) of Lemma \ref{lemma-properties-form}.

\medskip\noindent
Condition (3) is part (3) of Lemma \ref{lemma-properties-form}.

\medskip\noindent
Condition (4) is part (2) of Lemma \ref{lemma-properties-form}.

\medskip\noindent
Condition (5) follows from the fact that $(C_i \cdot C_j)$ is
the degree of an invertible module on $C_i$ which is divisible
by $[\kappa_i : k]$, see Varieties, Lemma \ref{varieties-lemma-divisible}.

\medskip\noindent
The genus formula proved in Lemma \ref{lemma-genus-formula}
tells us that the numerical type has the genus as stated, see
Definition \ref{definition-genus}.
\end{proof}

\begin{definition}
\label{definition-numerical-type-model}
In Situation \ref{situation-regular-model} the
{\it numerical type associated to $X$} is the numerical
type described in Lemma \ref{lemma-numerical-type-of-model}.
\end{definition}

\noindent
Now we match minimality of the model with minimality of the type.

\begin{lemma}
\label{lemma-numerical-type-minimal-model}
In Situation \ref{situation-regular-model}. The following
are equivalent
\begin{enumerate}
\item $X$ is a minimal model, and
\item the numerical type associated to $X$ is minimal.
\end{enumerate}
\end{lemma}

\begin{proof}
If the numerical type is minimal, then there is no $i$ with
$g_i = 0$ and $(C_i \cdot C_i) = -[\kappa_i: k]$, see
Definition \ref{definition-type-minimal}.
Certainly, this implies that none of the curves $C_i$
are exceptional curves of the first kind.

\medskip\noindent
Conversely, suppose that the numerical type is not minimal.
Then there exists an $i$ such that $g_i = 0$ and
$(C_i \cdot C_i) = -[\kappa_i: k]$.
We claim this implies that $C_i$ is an exceptional curve
of the first kind. Namely, the invertible sheaf
$\mathcal{O}_X(-C_i)|_{C_i}$ has degree $-(C_i \cdot C_i) = [\kappa_i : k]$
when $C_i$ is viewed as a proper curve over $k$, hence
has degree $1$ when $C_i$ is viewed as a proper curve over $\kappa_i$.
Applying
Algebraic Curves, Proposition \ref{curves-proposition-projective-line}
we conclude that $C_i \cong \mathbf{P}^1_{\kappa_i}$ as schemes
over $\kappa_i$. Since the Picard group of $\mathbf{P}^1$
over a field is $\mathbf{Z}$, we see that the normal sheaf
of $C_i$ in $X$ is isomorphic to $\mathcal{O}_{\mathbf{P}_{\kappa_i}}(-1)$
and the proof is complete.
\end{proof}

\begin{lemma}
\label{lemma-nonuniqueness}
Let $C$ be a smooth projective curve over $K$ with $H^0(C, \mathcal{O}_C) = K$
and genus $0$. If there is more than one minimal model for $C$, then
the special fibre of every minimal model is isomorphic to $\mathbf{P}^1_k$.
\end{lemma}

\begin{proof}
Let $X$ be some minimal model of $C$. The numerical type associated to $X$
has genus $0$ and is minimal (Definition \ref{definition-numerical-type-model}
and Lemma \ref{lemma-numerical-type-minimal-model}).
Hence by Lemma \ref{lemma-genus-zero} we see that
$X_k$ is reduced, irreducible, has $H^0(X_k, \mathcal{O}_{X_k}) = k$,
and has genus $0$. Let $Y$ be a second minimal model for $C$
which is not isomorphic to $X$. By
Resolution of Surfaces, Lemma \ref{resolve-lemma-birational-regular-surfaces}
there exists a diagram of $S$-morphisms
$$
X = X_0 \leftarrow X_1 \leftarrow \ldots \leftarrow X_n = Y_m
\to \ldots \to Y_1 \to Y_0 = Y
$$
where each morphism is a blowup in a closed point. We will prove the
lemma by induction on $m$. The base case is $m = 0$; it is
true in this case because we assumed that $Y$ is minimal
hence this would mean $n = 0$, but $X$ is not isomorphic
to $Y$, so this does not happen, i.e., there is nothing to check.

\medskip\noindent
Before we continue, note that $n + 1 = m + 1$ is equal to the number of
irreducible components of the special fibre of $X_n = Y_m$ because
both $X_k$ and $Y_k$ are irreducible. Another observation we
will use below is that if $X' \to X''$ is a morphism
of regular proper models for $C$, then $X' \to X''$ is an isomorphism
over an open set of $X''$ whose complement is a finite set
of closed points of the special fibre of $X''$, see
Varieties, Lemma
\ref{varieties-lemma-modification-normal-iso-over-codimension-1}.
In fact, any such $X' \to X''$ is a sequence
of blowing ups in closed points (Resolution of Surfaces, Lemma
\ref{resolve-lemma-proper-birational-regular-surfaces}) and the
number of blowups is the difference in the number of irreducible
components of the special fibres of $X'$ and $X''$.

\medskip\noindent
Let $E_i \subset Y_i$, $m \geq i \geq 1$ be the curve which is contracted
by the morphism $Y_i \to Y_{i - 1}$. Let $i$ be the biggest index such
that $E_i$ has multiplicity $> 1$ in the special fibre of $Y_i$.
Then the further blowups $Y_m \to \ldots \to Y_{i + 1} \to Y_i$
are isomorphisms over $E_i$ since otherwise $E_j$ for some $j > i$
would have multiplicity $> 1$. Let $E \subset Y_m$ be the inverse
image of $E_i$. By what we just said $E \subset Y_m$ is an
exceptional curve of the first kind. Let $Y_m \to Y'$ be the
contraction of $E$ (which exists by Resolution of Surfaces, Lemma
\ref{resolve-lemma-contract-when-quasi-projective}). The morphism
$Y_m \to X$ has to contract $E$, because $X_k$ is reduced.
Hence there are morphisms $Y' \to Y$ and $Y' \to X$ (by
Resolution of Surfaces, Lemma
\ref{resolve-lemma-factor-through-contraction})
which are compositions of at most $n - 1 = m - 1$ contractions
of exceptional curves (see discussion above). We win by induction on $m$.
Upshot: we may assume that the special fibres of all of
the curves $X_i$ and $Y_i$ are reduced.

\medskip\noindent
Since the fibres of $X_i$ and $Y_i$ are reduced, it has to be the
case that the blowups $X_i \to X_{i - 1}$ and $Y_i \to Y_{i - 1}$
happen in closed points which are regular points of the special fibres.
Namely, if $X''$ is a regular model for $C$ and if $x \in X''$
is a closed point of the special fibre, and
$\pi \in \mathfrak m_x^2$, then the exceptional fibre $E$ of the
blowup $X' \to X''$ at $x$ has multiplicity at least $2$ in the
special fibre of $X'$ (local computation omitted).
Hence $\mathcal{O}_{X''_k, x} = \mathcal{O}_{X'', x}/\pi$ is
regular (Algebra, Lemma \ref{algebra-lemma-regular-ring-CM}) as claimed.
In particular $x$ is a Cartier divisor on the unique
irreducible component $Z'$ of $X''_k$ it lies on
(Varieties, Lemma \ref{varieties-lemma-regular-point-on-curve}).
It follows that the strict transform $Z \subset X'$ of $Z'$
maps isomorphically to $Z'$
(use Divisors, Lemmas \ref{divisors-lemma-strict-transform} and
\ref{divisors-lemma-blow-up-effective-Cartier-divisor}).
In other words, if an irreducible component $Z$ of $X_i$
is not contracted under the map $X_i \to X_j$ ($i > j$)
then it maps isomorphically to its image.

\medskip\noindent
Now we are ready to prove the lemma.
Let $E \subset Y_m$ be the exceptional curve of the first kind
which is contracted by the morphism $Y_m \to Y_{m - 1}$. If $E$ is
contracted by the morphism $Y_m = X_n \to X$, then there is a factorization
$Y_{m - 1} \to X$ (Resolution of Surfaces, Lemma
\ref{resolve-lemma-factor-through-contraction})
and moreover $Y_{m - 1} \to X$ is a sequence of blowups
in closed points (Resolution of Surfaces, Lemma
\ref{resolve-lemma-proper-birational-regular-surfaces}).
In this case we lower $m$ and we win by induction.
Finally, assume that $E$ is not contracted by the morphism $Y_m \to X$.
Then $E \to X_k$ is surjective as $X_k$ is irreducible
and by the above this means it is an isomorphism.
Hence $X_k$ is isomorphic to a projective line as desired.
\end{proof}

\bibliographystyle{plain}
\bibliography{my}
\end{document}
